\section{The actual model, its points, and the symplectic trace inequality}
\label{sec:symplectic-model-comparison}

\status{This chapter reconstructs the passage from the actual
\texttt{Model} to the weighted intersection theorem, including the
point, tangent, trace and orbit comparisons.  The model is the one
constructed in MX, on the unchanged arithmetic datum of AX.  The
arithmetic assertions \texttt{Fin0}, \texttt{NoEnd} and
\texttt{AllIrr RigidAt} are stated explicitly below and are not
deduced from a dimension count.  Together with these assertions,
the already reconstructed rank-one theorems and the stated
classical local duality theorem supply the full comparison.
Neither an actual automorphic image theorem nor a finite
automorphic catalogue is inferred from this chapter.}

\subsection{The objects that must agree}

Fix an odd prime $\ell$, the original coefficient DVR
$O\subset\overline{\mathbb Q}_\ell$, its fraction field $E$,
and a datum $D$ satisfying \texttt{Good}.  Put
$K=\overline{\mathbb Q}_\ell$ and $G=G_{\mathbb Q}$ on the
specified algebraic closure.  All inertia and decomposition
groups retain the places chosen in $D$.  The four diagonal
characters and their two auxiliary weights are
\[
 (\theta_0,\theta_1,\theta_2,\theta_3)
       =(\alpha\Psi,\Psi,\alpha,1),\qquad
 (\lambda_i)=(1,0,1,0),\quad(\mu_i)=(1,1,0,0).
 \tag{PX.1}
\]
Thus $w_{ij}=(\lambda_i-\lambda_j,\mu_i-\mu_j)$, and
\[
 w_\alpha=(1,0),\quad w_\beta=(-1,1),\quad
 w_\Psi=(0,1),\quad w_\theta=(1,1).
 \tag{PX.2}
\]
The torus $(t_1,t_2)$ acts on an entry of weight $w=(a,b)$
by $t_1^at_2^b$, equivalently by conjugation with
$\operatorname{diag}(t_1t_2,t_2,t_1,1)$.

Let $M$ be an actual \texttt{Model} $D$.  Its weighted
presentation is $W:O[[z_1,\ldots,z_d]]\twoheadrightarrow R$,
with homogeneous kernel, origin $x_R:R\to O$ and
$\mathbb Z^2$-pieces $R_w$.  Write
\[
 B=\bigoplus_{w\in\mathbb Z^2}R_w\ \subset R,\qquad
 A=\bigoplus_{n\in\mathbb Z}R_{(0,n)}\ \subset B.
 \tag{PX.3}
\]
These are the finite-sum graded subalgebras; a completed ring
is not silently being declared a direct sum.  The model also
has a quotient $\pi:R\twoheadrightarrow S$ whose kernel is
generated by
\[
 F_0,\ldots,F_r\in R_{-w_\Psi},\qquad
 E_1,\ldots,E_{r'}\in R_{w_\Psi},\qquad
 Z_1,\ldots,Z_m\in R_0\cap\ker x_R,
 \quad r\ge1.
 \tag{PX.4}
\]
Its universal representation $\rho^u:G\to\operatorname{GSp}_4(S)$
is an actual \texttt{IsPt} lift, is generated by its matrix
entries in the source sense, and has homogeneous entry lifts
of weight $w_{ij}$ and origin $\theta_i\delta_{ij}$.
Let
\[
 \mathcal S=\bigoplus_w\pi(R_w)\subset S,
 \qquad \varphi:A\longrightarrow\mathcal S,
 \qquad I=(F_0,\ldots,F_r,E_1,\ldots,E_{r'},Z_1,\ldots,Z_m)_A.
 \tag{PX.5}
\]
We denote the induced grading of $A$ by $A_n$ and its origin
by $x_0$.  A point $\xi:\mathcal S\to K$ gives
$\rho_\xi(g)_{ij}=\xi(\rho^u(g)_{ij})$.

Here are the three retained arithmetic assertions, with their
literal quantifiers.  \texttt{Fin0} says that for every complete
Noetherian local $O$-algebra with the prescribed residue field,
and every actual \texttt{IsPt} lift over it, the $O$-algebra
generated by all entry monomials of total weight $(0,0)$ is a
finite $O$-module.  \texttt{NoEnd} says that an actual
algebraic-valued point which is block diagonal for the partition
$\{0,2\}\sqcup\{1,3\}$ cannot have its $\{0,2\}$ block
span all $2$-by-$2$ matrices.  \texttt{AllIrr RigidAt} says that
for every actual point $\rho$ whose matrix span is $M_4(K)$,
every actual dual-number point reducing to $\rho$ has derivative
$[X,\rho]$ for some $X\in\mathfrak{sp}_4(K)$.
These assertions concern all the points on this same datum,
not a selected finite family or a generic fibre alone.

We use the finite and rational local Tate duality and Euler
formulas already stated in LD/LK and DX, including their exact
transport to the chosen decomposition groups.  All remaining
commutative-algebra citation leaves are Noetherian formal
power series, Nakayama, exact completion on finite modules,
and the weak Nullstellensatz.  Precise references are Stacks
Project tags \href{https://stacks.math.columbia.edu/tag/0306}{0306},
\href{https://stacks.math.columbia.edu/tag/00DV}{00DV},
\href{https://stacks.math.columbia.edu/tag/00MA}{00MA}, and
\href{https://stacks.math.columbia.edu/tag/00FV}{00FV}.
The substitutions and invariant-theory arguments needed here
are proved below rather than left as additional suppliers.

\paragraph{PX1: actual gradings and generation in degrees at most two.}
\label{claim:PX1}
The sums in (PX.3) are direct; $A=\bigoplus_nA_n$ is a
graded $O$-algebra.  Its degree-zero ring is complete
Noetherian local with the original residue field.  It is
generated over $A_0$ by finitely many homogeneous elements
of degrees $-2,-1,1,2$, and $x_0(A_n)=0$ for $n\ne0$.
This is the literal \texttt{aGraded}/\texttt{aW} conclusion.

\begin{proof}
For a formal series let $p_w$ retain exactly the monomials
of weight $w$.  The kernel of $W$ is stable under every
$p_w$.  If a finite sum of elements of different weights
vanishes in $R$, lift the summands and apply each $p_w$;
each summand vanishes.  This proves both direct-sum assertions.
Products add weights because the coefficient of every monomial
in a formal product is a finite sum.

We record a useful finite-monomial argument.  For any fixed
weight $w$, the set of exponent vectors in $\mathbb N^d$ of
weight $w$ has finitely many coordinatewise minimal elements.
Here the elementary finiteness statement follows by induction
on $d$: an infinite sequence has an infinite subsequence with
nondecreasing first coordinate, and repeat on the remaining
coordinates; an antichain is therefore finite.  Every exponent
of weight $w$ dominates one of these minimal exponents, with
difference of weight zero.  Assign each monomial to the first
such minimal exponent.  Consequently every weight-$w$ series
is a finite sum of these monomials times weight-zero series.
Similarly, the nonzero minimal exponent vectors of weight zero
generate that monoid: subtract one and induct on total degree.
There are only finitely many such generators.  Thus the full
weight-zero ring is a quotient of a formal power series ring
over $O$ on these nonconstant monomials.  Substitution is
well-defined because their degrees are positive.  A series
with unit constant coefficient has an inverse of weight zero,
obtained either by a geometric series or by projecting an
inverse.  This proves locality, completeness, Noetherianness
and the original residue field.  The same argument applies
after a homogeneous quotient.

Each original variable has first weight in $\{-1,0,1\}$
and second weight between $-1$ and $1$.  A monomial of first
weight zero is a product of variables of first weight zero
and pairs consisting of one variable of first weight $1$
and one of first weight $-1$.  Remove one such factor and
induct on degree.  Apply the preceding finite-monomial
decomposition to each full weight $(0,n)$.  It shows that
$A$ is generated over $A_0$ by the finitely many corresponding
single variables and pairs.  Their second weights lie between
$-2$ and $2$; those of weight zero are absorbed into $A_0$.
Finally a nonzero-weight series has constant coefficient zero,
which is exactly the origin assertion.  Notice that none of
these steps replaces a formal series by an unjustified
infinite sum in the algebraic direct sum.
\end{proof}

\paragraph{PX2: the invariant chart and its exact dimensions.}
\label{claim:PX2}
The source \texttt{SmoothPt} condition on $B$ induces
$\texttt{SmoothPt}(A_n,x_0,\mathtt{mDims}(r,0,r',0,m))$.
Explicitly the saturated cotangent ranks of $A$ are
\[
 \begin{array}{c|rrrrr}
 n&-2&-1&0&1&2\\\hline
 \operatorname{rank}_O C_n(A,x_0)&0&r&m&r'+1&1.
 \end{array}
 \tag{PX.6}
\]
The chart and the saturation in this statement are over $E$;
no integral formal smoothness is asserted.

\begin{proof}
Write $J=\ker(B\to O)$ and $J_A=\ker(A\to O)$.  The source
chart $\Phi:B\to E[[Y_i:i\in T]]$ respects both weights
and the origin.  Its linear coefficients identify
\[
 (J/J^2)\otimes_O E\ \simeq\ E^T.
 \tag{PX.7}
\]
Indeed its \texttt{inj} field says precisely that a zero
linear part lies in the $O$-saturation of $J^2$, and its
\texttt{onto} field supplies a diagonal basis of nonzero
linear parts.  Projection to the appropriate weight and
subtraction of the constant term make these basis lifts
homogeneous.  This also proves directly that (PX.7) respects
weights.

For clarity we justify the use of formal coordinates here.
The finite-monomial proof of PX1 makes $B$ Noetherian: it is
generated by the finitely many $z_i$ over the complete
Noetherian full weight-zero ring.  Complete $B\otimes_O E$
at the augmentation ideal.  Homogeneous lifts of the basis
in (PX.7) give a surjection from $E[[Y_i]]$ to this
completion, by successive approximation and Nakayama.
Composing with $\Phi$ gives a formal substitution with
invertible linear matrix.  Such a substitution is invertible:
after fixing the inverse through degree $q$, solve its
degree-$q+1$ error with the inverse linear matrix.  The
successive corrections converge in the variable-adic
topology.  It follows that the surjection and the completed
chart are isomorphisms.  The construction is equivariant
for both weights, since every correction may be projected
to its prescribed weight.

Passing to first-weight zero commutes with these completions
in this situation.  Here is the needed topology check,
which prevents an unwarranted use of formal invariants.
A first-weight-zero monomial in the original variables
splits into singles and opposite-weight pairs.  Any such
monomial of total degree at least $2q$ has at least $q$
nonconstant blocks, and lies in $J_A^q$.  Conversely
$J_A^q\subset J^q$.  On a homogeneous component an element
of $J^{2q}$ has a formal representative of order at least
$2q$: lift products of $2q$ augmentation elements, then
project to that weight.  Grouping its monomials by the
finitely many possible first $q$ blocks proves the same
inclusion for series.  Taking the finite sum of components
gives
\[
        J^{2q}\cap A\subset J_A^q\subset J^q\cap A.
 \tag{PX.8}
\]
The argument survives localization in $O\setminus\{0\}$.
Weight-zero projection is an additive retraction compatible
with each of these filtrations.  Thus compatible invariant
truncations lift compatibly, and (PX.8) identifies the
completion of $A\otimes E$ with the invariant part of the
completed $B\otimes E$.  This is also an explicit proof of
the required exactness of taking this torus's invariants.

The source dimension table has one variable $Y_\alpha$
of weight $(1,0)$, one $Y_\beta$ of weight $(-1,1)$,
one $Y_\theta$ of weight $(1,1)$, $r'$ variables of
weight $(0,1)$, $r$ of weight $(0,-1)$, and $m$ of
weight zero.  There is exactly one negative first-weight
variable, namely $Y_\beta$.  A monomial of first weight
zero is uniquely a monomial in
\[
 Y_i\quad(w_i^{(1)}=0),\qquad
 U_\alpha=Y_\alpha Y_\beta,\qquad
 U_\theta=Y_\theta Y_\beta.
 \tag{PX.9}
\]
The exponent of $Y_\beta$ is the sum of the exponents of
the two positive-first-weight variables.  This gives an
isomorphism of formal rings, coefficient by coefficient;
there are no relations among the displayed generators.
Their second weights are respectively unchanged, $1$, and
$2$.  Restrict the source chart and relabel these monomials.
The completed-ring identification proves its linear map is
an isomorphism on the rational cotangent space.  Clearing
the finitely many denominators supplies the source
\texttt{onto} witnesses in $A$, and zero linear part is
equivalent to membership in the saturation of $J_A^2$.
This proves every field of \texttt{IsChart} and (PX.6).
\end{proof}

\paragraph{PX3: global tangent vanishing and all three spanning conditions.}
\label{claim:PX3}
Using RV7 and the ordinary tangent consequences in DX, the
equations (PX.4) satisfy
\[
 \mathtt{SpansCot}(-1,F),\qquad
 \mathtt{SpansCot}(-2,\varnothing),\qquad
 \mathtt{SpansCot}(0,Z).
 \tag{PX.10}
\]
Together with PX2 and $r\ge1$, this is the literal condition
$\mathtt{M1}$ for $M.\mathtt{eqns}$.

\begin{proof}
We first prove the actual global tangent assertion used in
the spanning test.  Suppose $\rho':G\to M_4(O[\epsilon])$
is an actual point, reduces to $\rho_0$, and its derivative
has only weight $w=0$ or $w=-w_\Psi$.  Set
\[
 z(g)_{ij}=\epsilon^{-1}(\rho'(g)-\rho_0(g))_{ij}
                                  \theta_j(g)^{-1}.
 \tag{PX.11}
\]
Expanding the representation identity shows that each entry
is a cocycle for $\theta_i\theta_j^{-1}$; the fixed
similitude identity places $z$ in $\mathfrak{sp}_4$.
The entries are continuous.  One direct check is that a
map from a complete local presentation to $O[\epsilon]$
sends the first component of its maximal ideal into
$\mathfrak m_O$, and the second component of its
$(q+1)$st power into $\mathfrak m_O^q$.  Thus the
given adic continuity implies continuity of the derivative.

For $w=0$ the entries are additive diagonal characters.
DX's weight-zero ordinary tangent calculation makes them
zero on inertia at $\ell$.  At $p\in S_b$ they vanish on
$I_p\cap\ker\psi$, and away from $T$ on
$I_p\cap\ker\bar\rho$.  The omitted subgroups have finite
index.  An additive characteristic-zero cocycle zero on a
finite-index subgroup is zero on all inertia, by taking a
power.  Transport by conjugation handles every place.
A continuous additive character of $G_{\mathbb Q}$
unramified everywhere is zero: each finite quotient of its
compact image cuts out an everywhere-unramified abelian
extension of $\mathbb Q$, which is trivial by the class
number of $\mathbb Q$ (equivalently Minkowski and class
field theory).  Passing through all finite quotients gives
the assertion for the original continuous character.

For $w=-w_\Psi$, only the entries $(2,0)$ and $(3,1)$
can occur; the symplectic identity relates them.  The first
entry is a cocycle for $\Psi^{-1}=\psi^{-1}\chi_{\rm cyc}^{-e}$.
DX's negative ordinary tangent calculation makes it a
coboundary on $G_\ell$ after passing to $E$.  At all other
inertia groups the same finite-index argument makes it a
coboundary.  Explicitly, if a cocycle vanishes on a
finite-index subgroup of inertia, it vanishes on the
kernel of its character by taking powers; choosing one
element with character different from $1$ then writes it
as $(\gamma-1)c$.  If the character is trivial it is
already zero.  RV7 applies to precisely this cocycle,
including its local condition at $\ell$, and gives
$z(g)_{20}=(\Psi(g)^{-1}-1)c$ globally.  The derivative
of the normalization $h(\rho')_{20}=0$ at $g_0$ gives
\[
             c(\alpha(g_0)-1)(\Psi(g_0)-1)=0.
 \tag{PX.12}
\]
Both factors are units by \texttt{Good}.  Hence $c=0$,
and the symplectic relation kills the other entry.  This
proves the full \texttt{tanGlob} assertion in both weights.

Now work first on $B$.  Its cotangent lattice of any
fixed weight is finite: a homogeneous series modulo the
square of the augmentation is a combination of the
finitely many original variables of that weight.  Suppose
the $F_i$ failed to span the rational cotangent space at
$-w_\Psi$.  A nonzero $E$-linear functional annihilating
them can be multiplied by one denominator to take values
in $O$ on this finite lattice.  It defines an $O$-valued
point derivation of the specified weight.  The same
argument applies to the $Z_i$ at weight zero.
Such a derivation on $B$ extends to $R$: on the formal
presentation use the finite linear part at the origin;
homogeneous kernel relations are killed, and the
quadratic terms give the Leibniz rule.  Its weight forces
it to kill the other two groups of equations in (PX.4).
Consequently $x_R+\epsilon\delta$ factors through $S$.
The universal representation gives an actual tangent
point of the kind just proved zero.  The model's
entry-generation property then makes the two maps
$S\to O[\epsilon]$ equal, so $\delta=0$, a contradiction.
Thus $F$ and $Z$ span on $B$ after rationalization.

Inclusion $A\subset B$ induces surjections of cotangent
spaces at weights $-1\mapsto-w_\Psi$ and $0\mapsto0$:
the homogeneous pieces themselves agree.  PX2 gives
their ranks at most $r,m$, while the original model
chart gives the target ranks exactly $r,m$.
They are therefore rational isomorphisms, transferring
the just-proved spanning statements to $A$.  The
degree-$-2$ cotangent has rational rank zero by PX2;
its defining saturation makes it torsion free, so it
vanishes and is spanned by the empty family.  Finally,
spanning after rationalization is exactly the displayed
\texttt{SpansCot} predicate: for each integral class
clear its denominator, without asserting integral
generation or changing the integral cotangent quotient.
\end{proof}

\paragraph{PX4: the exact equation ideal and lifting every algebraic point.}
\label{claim:PX4}
The map $\varphi$ has kernel $I$, is onto each
first-weight-zero piece of $\mathcal S$, and preserves
its second weight.  Every $x:A\to K$ with $x(I)=0$
extends to $\xi:\mathcal S\to K$ with $x=\xi\varphi$.
These are the literal \texttt{kerA} and \texttt{ptLift}
statements; no irreducibility hypothesis is used.

\begin{proof}
The ideal in (PX.4) is homogeneous, so $S$ has the
induced weighted presentation and its finite homogeneous
sum $\mathcal S$ has a genuine bigrading by the projection
argument of PX1.  If $f\in A$ lies in $\ker\pi$, write
it as a finite linear combination of the generators in
(PX.4) with coefficients in $R$.  For each of the
finitely many full weights of $f$, project this identity
to that weight.  Since every generator has first weight
zero, only the first-weight-zero components of its
coefficient survive, and these lie in $A$.  The resulting
finite identity puts $f$ in $I$.  The reverse inclusion
is immediate.  A homogeneous lift of a class in
$\pi(R_{(0,n)})$ proves the other two assertions about
$\varphi$.

Let $C=\varphi(A)=\mathcal S^{\mathbb G_m^{(1)}}$.
The point $x$ descends to $C$.  Projection
$P:\mathcal S\to C$ to first weight zero is $C$-linear
and has $P(1)=1$.  Hence the fibre
$\mathcal S\otimes_{C,x}K$ is nonzero: $xP$ induces
a $K$-linear map from this fibre to $K$ sending $1$ to
$1$.  It is of finite type over $K$.  Indeed the
finite-monomial decomposition in PX1 says that the full
graded algebra is generated over its full weight-zero
piece by the finitely many original variables; that
piece is contained in $C$.  Choose a maximal ideal of
the nonzero fibre.  The weak Nullstellensatz makes its
residue field $K$, yielding the desired extension.
Thus surjectivity on points is a proved Reynolds
projection and finite-type argument, not an assumption
that every point of a subalgebra extends automatically.
\end{proof}

\paragraph{PX5: algebraic entry generation, actual points and condition M2.}
\label{claim:PX5}
Assume \texttt{Fin0}.  Each homogeneous piece of
$\mathcal S$ is the $O$-linear span of the entry
monomials of that full weight.  Two algebra maps out
of $\mathcal S$ into any $O$-algebra agree if their
evaluated representations agree.  Every point
$\xi:\mathcal S\to K$ can be rescaled by the full
two-dimensional torus to the restriction of a map
$S\to K$, so its representation becomes an actual
$D.\mathtt{EPt}$.  Finally the degree-zero image of
$A/I$ is finite over $O$, the literal condition
$\mathtt{M2}$.

\begin{proof}
Let $C_0$ be the $O$-algebra generated by the entry
monomials of weight zero.  It is finite by
\texttt{Fin0}, applied to the actual universal lift.
Choose finitely many homogeneous entry differences
(subtract a scalar from diagonal entries) whose
relative tangent classes generate that of $S$.
They exist by the model's generation assertion: a
linear functional killing all these classes would
give two different maps into the residue-field dual
numbers with identical representations.  Nakayama
and successive approximation show that these finitely
many differences generate $S$ topologically.
The resulting new presentation has homogeneous kernel
for the original grading.  To see this directly, lift
each chosen difference to a homogeneous series in the
original $O[[z_s]]$.  Its constant coefficient can be
taken zero: the origin descends to $S$, since every
equation in (PX.4) has origin zero, and the differences
were centered at that origin.  Substitute these lifts
into a series in the new variables.  Their positive
formal order makes every output coefficient a finite
sum, so substitution commutes with each weight
projection.  The original presentation kernel is
homogeneous, hence so is the new kernel.  The same
identity shows that every old homogeneous element
has a new homogeneous preimage.  No assertion that
an arbitrary image of a weight piece is closed is
needed here.

The finite ring $C_0\subset S$ is complete over $O$.
It is local: the usual finite-ring decomposition over
a complete local ring would otherwise give a nontrivial
idempotent inside the local ring $S$.  Its maximal
ideal is $C_0\cap\mathfrak m_S$, and its topology is
equivalent to its $\mathfrak m_O$-adic topology.
The weight-zero monomials in the chosen differences
have positive order in $S$ and lie in that maximal
ideal.  Their formal series therefore converge in
$C_0$, and the inclusion in $S$ identifies their
limits with the original series.  For any weight $w$,
the finite-monomial argument of PX1 writes a
homogeneous series as a finite sum of entry monomials
times such weight-zero series.  Thus its value is
an $O$-linear combination of entry monomials of
weight $w$.  This proves the exact span assertion,
including the removal of a possible topological
closure.  It immediately proves uniqueness of an
algebra map from all its matrix entries.  Torus
rescaling defines an algebra map on $\mathcal S$,
because weights add.

Here is the point-extension argument with its
convergence condition included.  For presentation
variables $z_s$ put $a_s=v(\xi(z_s))$ when their
values are nonzero, and write $w_s\in\mathbb Z^2$
for their weights.  If $n_s\ge0$ are integers, not
all zero, with $\sum n_sw_s=0$, then
\[
                   \sum n_sa_s>0.
 \tag{PX.13}
\]
To prove this, the corresponding monomial $u$ has
weight zero and positive formal order.  If
$\xi(u)^{-1}$ were integral over $O$, a monic
polynomial annihilating it could be reversed to
give a polynomial in $u$ with constant coefficient
$1$, hence a unit in the weight-zero formal ring,
whose image under $\xi$ is zero.  This is impossible.
The valuation characterization of integrality in
an algebraic extension of a complete DVR proves
(PX.13).

The finite strict linear inequalities
$a_s+\langle h,w_s\rangle>0$ therefore have a
solution $h\in\mathbb Q^2$.  For completeness,
failure of these inequalities separates the affine
subspace $a+\{(\langle h,w_s\rangle)_s\}$ from
the positive orthant, producing a nonzero vector
$q_s\ge0$ with $\sum q_sw_s=0$ and
$\sum q_sa_s\le0$.  All data are rational, so
a rational such vector can be chosen by solving
the finite linear equalities on a supporting face.
Clearing denominators contradicts (PX.13).
A real feasible vector may likewise be approximated
by a rational one, because the inequalities are
strict.  Choose $t_1,t_2\in K^\times$ of valuations
the two coordinates of $h$; roots of a uniformizer
provide them.  All scaled variable values now have
positive valuation.  Variables whose values were
zero impose no condition.

The scaled values lie in one finite extension of
$E$.  Adjoin to $O$ these finitely many values,
generators of the finite module $\xi(C_0)$, and
$y_s^{N_s}/\varpi$ with $N_s$ large enough that
these last elements are integral.  This gives a
finite $O$-algebra $C\subset K$ in which all $y_s$
are in the radical of $\varpi C$.  It is complete,
so the power series may be evaluated in $C$.
This evaluation is initially defined on the formal
source, before descent to $S$.  Let $P_0$ be the
full-weight-zero subring of that source.  Its map
$W_0:P_0\to C_0=\mathcal S_0$ is local and continuous:
both rings are complete local with residue field $k$,
and every nonconstant monomial maps into
$C_0\cap\mathfrak m_S$.  The restriction $\xi|_{C_0}$
is continuous because it is $O$-linear on a finite
$O$-module.  Evaluation on $P_0$ is continuous as
well, since its finitely many nonconstant monomial
generators have positive valuation.  It agrees with
$\xi W_0$ on these generators, because their full
weight is zero.  Successive approximation therefore
proves agreement on all of $P_0$.
For every other weight the finite-monomial
decomposition gives agreement with $\xi W$ after
the indicated torus factor.  Consequently evaluation
kills each homogeneous component of a relation.
For each $n$, sum the complete weight components
whose weights occur in total degree less than $n$;
there are finitely many, and the remaining series
has order at least $n$.  Evaluation is continuous,
so letting $n$ grow shows that it kills the whole
relation.  Thus evaluation factors through $S$
without assuming that descent in advance.  The
resulting representation is an actual point by its
displayed \texttt{IsPt} witness $\rho^u$.

Finally PX4 identifies $A/I$ with a subalgebra of
$\mathcal S$.  Its degree-zero piece maps into the
finite module $C_0$ by the span assertion.  A
submodule of a finite module over the Noetherian
ring $O$ is finite.  This proves $\mathtt{M2}$.
\end{proof}

\paragraph{PX6: normalized points and the three possible matrix patterns.}
\label{claim:PX6}
Every actual point has continuous entries in the integers
of one finite extension of $E$, strictly congruent to
$\rho_0$ in the $\ell$-adic norm.  Every diagonal matrix
unit $e_{ii}$ belongs to the span of its matrices, already
with integral coefficients in a common finite extension.
For an actual point $\rho$, write $i\rightsquigarrow j$
if $\rho(g)_{ij}\ne0$ for some $g$.  Exactly the following
exhaustive alternatives are needed:
\[
 \begin{split}
 &\text{the matrix span of $\rho$ is $M_4(K)$;}\\
 &i\ne j,\ i\rightsquigarrow j\ \Longrightarrow\
                                      j\not\rightsquigarrow i;\\
 &0\rightsquigarrow2,\ 2\rightsquigarrow0,
 \quad\text{and all cross-block entries vanish in one direction.}
 \end{split}
 \tag{PX.14}
\]
Here cross blocks refer to $\lambda=1$ and $\lambda=0$.
The second alternative is the source \texttt{DiagType}.
This conclusion uses the residual vanishing already in
\texttt{Good} and finite local Tate duality and Euler at
the primes in $S_b$; no automorphy or irreducibility
criterion is an unstated input.

\begin{proof}
An algebra map $y:C\to K$ from an actual complete local
coefficient algebra is automatically bounded in the
needed sense.  If $c\in\mathfrak m_C$ and $|y(c)|\ge1$,
then $y(c)^{-1}$ is integral over $O$.  Reverse a monic
polynomial annihilating it.  Its value at $c$ is a unit,
since its constant coefficient is $1$, but its image is
zero, a contradiction.  Thus $|y(c)|<1$ on
$\mathfrak m_C$.  Every element is a scalar from $O$
plus an element of that ideal, so all values have norm
at most $1$.  Finitely many generators of
$\mathfrak m_C$ give a uniform bound $a<1$ on its
image and a bound $a^q$ on its $q$th power.  This
proves continuity, including continuity of the
specialized representation.  A finite formal
presentation now puts all values in the finite
extension generated by its finitely many variable
values; convergence takes place in that complete
field.  The residual identity gives
\[
                  |\rho(g)_{ij}-\rho_0(g)_{ij}|<1.
 \tag{PX.15}
\]

We next spell out why normalization supplies the
matrix projectors.  AX3 and AX6 give that
$\rho(\phi)$ has order prime to $\ell$, is block
diagonal, and has different residual eigenvalues
on the two blocks.  A matrix of prime-to-$\ell$
order which is scalar modulo the maximal ideal
is that scalar's unique finite-order lift:
factor $X^d-1$ at that root, whose complementary
factor is a unit.  Applying this to each block
shows that the two block projectors are
polynomials in $\rho(\phi)$ with integral
coefficients.  In the $\{0,2\}$ block the
self-adjoint element
\[
 H=\rho(g_0)+\theta_0(g_0)\rho(g_0^{-1})
 \tag{PX.16}
\]
has zero off-diagonal entries by normalization.
Its diagonal difference is a unit, by AX6:
its residue is
$(\bar\alpha(g_0)-1)(\bar\Psi(g_0)-1)$.
Linear interpolation in this block gives $e_{00}$
and $e_{22}$.  The symplectic adjoint gives
$e_{33}$ and $e_{11}$.  The span of the group
matrices is closed under multiplication and
adjoint, because products are group matrices
and the adjoint of $\rho(g)$ is its prescribed
similitude times $\rho(g^{-1})$.  This proves
the projector assertion over the coefficient
algebra and after every specialization.

It follows that $\rightsquigarrow$ is reflexive
and transitive, and is carried by symplectic
duality to $j^*\rightsquigarrow i^*$, where
$i^*=3-i$.  Indeed a nonzero entry and the
two projectors supply the matrix unit $e_{ij}$;
multiply matrix units to prove transitivity.
The full matrix span is equivalent to every
pair being linked.

The substantive step is that a two-element
strongly connected component cannot straddle
the two $\lambda$ blocks.  Suppose it consists
of $s,u$ with $\lambda_s\ne\lambda_u$.
The component gives a two-dimensional
subquotient; products of its matrices have no
contribution through an intermediate vertex
outside this component.  The ordinary flag
of AX8 gives this local subquotient a stable
line at $\ell$.  In one of the two coordinate
charts that line has slope $v$ with $|v|\le1$;
interchange $s,u$ if necessary.  Both links
exist, so the chosen off-diagonal coefficient
is nonzero.  In a common coefficient DVR
$O'$, let $b\ne0$ generate the ideal of all
$\rho(g)_{su}$ and write
$\rho(g)_{su}=b c(g)$ with $c(g)\in O'$.
The reductions of the $c(g)$ are not all zero.
The matrix-product identity gives the residual
cocycle
\[
 z(g)=\frac{\overline{c(g)}}{\bar\theta_u(g)},
 \qquad \gamma=\bar\theta_s\bar\theta_u^{-1}.
 \tag{PX.17}
\]
It cannot be a coboundary.  If
$\overline{c(g)}=a(\bar\theta_s(g)-\bar\theta_u(g))$,
apply the integral linear combination of group
matrices giving $e_{ss}$.  Its normalized
off-diagonal entry gives $0=a$, contradicting
that some $\overline{c(g)}$ is nonzero.

On the other hand this cocycle has all the
local conditions of $D.\mathtt{XVan}(\theta_s/\theta_u)$.
Outside $T$, the actual unramified condition
kills it on inertia intersected with
$\ker\bar\rho$.  At $p\in S_b$ the local
$H^1$ of its residual character is zero.
Here is the precise duality calculation.
AX6 provides $\sigma_p\in G_p$ for which
$\gamma(\sigma_p)\ne1$ and the residual
cyclotomic value is $1$.  Thus both $H^0(k(\gamma))$
and $H^0(k(\gamma^{-1}\bar\chi_{\rm cyc}))$
vanish.  Local Tate duality gives $H^2=0$;
the local Euler formula at $p\ne\ell$ then
gives $\dim H^1=0$.  This is valid for
the actual finite coefficient extension $k$.

At $\ell$, the stable-line identity is
\[
 \rho(g)_{ss}v+b c(g)
             =(\rho(g)_{us}v+\rho(g)_{uu})v.
 \tag{PX.18}
\]
Apply it also to $\sigma_\ell$ and eliminate
$b$ and the linear slope term.  Reduction
gives
\[
 \overline{c(g)}(\bar\theta_u(\sigma_\ell)
                      -\bar\theta_s(\sigma_\ell))
 =\overline{c(\sigma_\ell)}
                     (\bar\theta_u(g)-\bar\theta_s(g)).
 \tag{PX.19}
\]
The discarded terms are multiples of
$v\rho(g)_{us}$ or $v\rho(\sigma_\ell)_{us}$;
their norms are strictly less than $1$ by
(PX.15).  The first parenthesis is nonzero
by AX6, so (PX.19) makes (PX.17) a local
coboundary.  The ratio $\theta_s/\theta_u$
is one of $\alpha,\beta,\theta$ or their
inverses, because the first weights differ.
The residual vanishing in \texttt{Good}
applies.  It remains valid after extension
of the residue field: expand the cocycle
and its local coboundary constants in a
basis over the original residue field and
apply the original vanishing coordinate by
coordinate.  Thus (PX.17) is a global
coboundary, the required contradiction.

Finish by considering the strongly connected
components on four vertices.  A three-vertex
component and its dual overlap, hence coincide;
this is impossible since duality has no fixed
vertex and a stable subset has even size.
A four-vertex component is the first case of
(PX.14).  If all components are singletons
we have \texttt{DiagType}.  Otherwise the
only permitted two-vertex component is one
of $\{0,2\}$ and $\{1,3\}$; duality gives
both.  Arrows between these two components
can point in at most one direction, or all
four vertices would be strongly connected.
This is exactly the third case in (PX.14).
\end{proof}

\paragraph{PX7: one-sided reducible points, two-sided irreducibility and M3.}
\label{claim:PX7}
Assume \texttt{Fin0} and \texttt{NoEnd}.  Every two-sided
point of $A/I$ gives, through PX4 and PX5, an irreducible
actual representation.  Moreover $A,I,x_0$ satisfy the
literal $\mathtt{M3}$ condition.  The latter conclusion
uses \texttt{Fin0} but does not need \texttt{NoEnd}.

\begin{proof}
First consider an actual representation of
\texttt{DiagType}.  Transitivity and the absence
of a two-way link make its diagonal entries
characters $\kappa_i$.  Restrict to $G_\ell$.
The ordinary flag of AX8 and the equality of
traces identify their multiset with the four
ordinary diagonal characters.  To see that
the matching is unique, restrict to $s_0$:
its four character values are distinct by
AX6.  Independence of distinct characters
then matches them as characters, not only
as values at one element.  Residual separation
at $\sigma_\ell$ preserves $\lambda$, and
symplectic duality preserves $i\mapsto3-i$.
There are precisely two resulting permutations:
\[
                 m=\mathrm{id}\quad\hbox{or}\quad
                 m=(2,3,0,1).
 \tag{PX.20}
\]
This last enumeration follows by choosing the
image of $0$ in $\{0,2\}$; duality fixes
the image of $3$, and the remaining two
images are then fixed.  Thus on $I_\ell$
we have $\kappa_i=\theta_{m(i)}$.

Each $\kappa_i$ is
$\chi_{\rm cyc}^{h_{m(i)}}\delta_i$, with
$(h_0,h_1,h_2,h_3)=(e+1,e,1,0)$ and
$\delta_i$ of finite order.  Here is the
global finiteness check.  Divide by the
displayed cyclotomic power.  On inertia
this quotient is trivial on the intersection
with one fixed open subgroup
$\ker\bar\rho\cap\ker\psi\cap\ker\omega$,
using the local conditions and AX8 at
$\ell$.  Raising to its index makes the
character unramified at every finite prime.
An everywhere-unramified abelian character
of $G_{\mathbb Q}$ is trivial, as in PX3.
Thus a fixed power of $\delta_i$ is trivial;
continuity gives its open kernel.  This
argument does not assume in advance that
an arbitrary algebraic-valued character
has finite image.

We claim that a link $i\rightsquigarrow j$
between distinct vertices forces $m(i)<m(j)$.
It is enough to prove this for a cover link,
with no vertex strictly between the two:
any link in a finite partially ordered
set is a chain of cover links.  For such
a link put
\[
 z(g)=\rho(g)_{ij}/\kappa_j(g),
 \qquad \eta=\kappa_i/\kappa_j.
 \tag{PX.21}
\]
There is no intermediate contribution to
matrix multiplication, so this is a
continuous $\eta$-cocycle.  If $m(i)>m(j)$,
then $\eta$ is a finite character times
a strictly negative power of cyclotomic,
since $e+1>e>1>0$.  The ordinary local
flag splits this reversed extension at
$\ell$: its higher-character line is
complementary to the lower-character
line, and the distinct $s_0$ eigenvalues
make the complement a graph.  Equivalently
the $2$-by-$2$ quotient has
$\rho(g)_{ij}=a(\kappa_i(g)-\kappa_j(g))$
on $G_\ell$.  At every other inertia
group the actual unramified equations
and a finite-index argument give the
same coboundary condition.  RV7 now
makes (PX.21) a global coboundary.
Applying the linear combination giving
$e_{ii}$ forces its constant to be zero,
so the entry vanishes for every $g$,
contrary to the link.  The claim follows.

When $m=\mathrm{id}$, every negative
second-weight entry therefore vanishes,
and $\rho(s_0)_{00}=\theta_0(s_0)$.
When $m=(2,3,0,1)$, every positive
second-weight entry vanishes, and
$\rho(s_0)_{00}=\theta_2(s_0)$.
These two distinguished values differ.
By PX5 a negative-weight entry must
occur in every entry monomial of
negative total second weight; similarly
for positive weight.  Hence these
entry vanishings imply respectively
$\mathtt{NegZero}$ and $\mathtt{PosZero}$
on the restricted point of $A$.
Conversely $\mathtt{PosZero}$ kills
every $(0,2)$ entry and
$\mathtt{NegZero}$ every $(2,0)$ entry,
because their weights are $(0,1)$ and
$(0,-1)$ and PX4 lifts those pieces.

Apply PX6 to a torus-scaled actual
representative of a two-sided point.
The diagonal-type case has just been
excluded.  In the block-triangular
case, scale by the first torus parameter
and let it tend to zero in the direction
which kills the cross-block entries.
This is an algebraic specialization,
not an analytic limiting assertion:
entry monomials have only nonnegative
powers of the parameter wherever their
values are nonzero, by the one-sided
block condition and PX5.  Evaluation
at zero is an algebra map on
$\mathcal S$ retaining its first-weight-zero
pieces.  PX5 rescales its resulting
representation to an actual point.
Its $(0,2)$ and $(2,0)$ links survive;
the projectors and these links make its
$\{0,2\}$ block span $M_2(K)$.
This contradicts \texttt{NoEnd}.  Only
the irreducible case remains.

For $\mathtt{M3}$ take a point $x_1$
with both signs zero, vanishing on $I$,
and satisfying the source's additional
inequality $x_1\ne x_0$.  Lift and
scale it.  The absent $(0,2)$ and
$(2,0)$ links exclude the first and
third cases of PX6, so its representation
has diagonal type.  Points with the
same degree-zero restriction have the
same diagonal entries: each is of
full weight zero.  They also share
diagonal type, since every product
$\rho(g)_{ij}\rho(h)_{ji}$ is a
degree-zero function and vanishes
for a diagonal-type representation.
They must have the same alternative
in (PX.20), because their $(0,0)$
entries at $s_0$ agree and the two
possible values are distinct.
Thus they all have the same sign
vanishing.  One of the two opposite
one-sided families in the definition
of $\mathtt{M3}$ is empty, exactly
as required.  The proof even gives
this conclusion without using the
extra inequality $x_1\ne x_0$.
\end{proof}

\paragraph{PX8: normalized conjugacy is torus conjugacy, with the correct orbit.}
\label{claim:PX8}
Two conjugate actual points, one of which is irreducible,
are conjugate by the torus in (PX.1).  Consequently, under
\texttt{Fin0} and \texttt{NoEnd}, conjugate representations assigned to
two-sided points of $A/I$ come from the same orbit of the
second one-dimensional torus on $A$.

\begin{proof}
Let $U\rho(g)=\rho'(g)U$ with $U$ invertible.
The normalization and (PX.15) force $U$
to be block diagonal.  Here is an elementary
way to see the separation argument, also
valid without diagonalizing either matrix.
In an off-block entry of
$U\rho(\phi)=\rho'(\phi)U$, subtract the
common residual diagonal matrix.  Choose
an off-block entry of $U$ of maximum norm.
Its coefficient is a difference of two
residually distinct diagonal values and
has norm $1$.  Every error term is an
off-block entry of norm at most that
maximum, multiplied by an error of
norm strictly less than $1$.
The ultrametric inequality gives a
strict contradiction unless the maximum
is zero.

Intertwining also holds for the two
elements $H,H'$ in (PX.16).  Their
$(0,2)$ and $(2,0)$ entries are zero,
and the required opposite diagonal
values are distinct by the same
unit separation from AX6 and (PX.15).
The $(0,2)$ and $(2,0)$ entries of
$UH=H'U$ therefore force
$U_{02}=U_{20}=0$.
Moreover $U^*U$ commutes with every
$\rho(g)$, since both representations
have the same similitude.  Irreducibility
means their matrix span is all $M_4(K)$,
so $U^*U$ is scalar.  This follows
directly by commuting with the matrix
units; no version of Schur's lemma
over an unspecified endomorphism
field is needed.  Thus $U$ is a
symplectic similitude.  Its block
form and the two zero corner entries
then force the remaining off-diagonal
entries $U_{13},U_{31}$ to vanish.
Its diagonal entries $d_i$ obey
$d_0d_3=d_1d_2$ and are nonzero.
Dividing by $d_3$ gives exactly
\[
       (d_0,d_1,d_2,d_3)/d_3
             =(t_1t_2,t_2,t_1,1),
 \qquad t_1=d_2/d_3,\ t_2=d_1/d_3.
 \tag{PX.22}
\]
Scalar matrices act trivially, proving
the claimed torus conjugacy.

Now let $\xi,\xi'$ lift two points of
$A/I$.  PX5 scales them by torus
elements $t,t'$ to actual points.
If these are conjugate, (PX.22)
applies because PX7 makes the two-sided
representations irreducible under the
two stated hypotheses.  Composition of scalings gives a
torus element carrying the evaluated
representation of $\xi$ to that
of $\xi'$.  Entry generation in
PX5 makes the corresponding algebra
maps equal on all of $\mathcal S$.
Restrict to $A$.  Its first weights
are zero, so the first parameter
disappears and only the second
parameter acts.  This is the exact
\texttt{SameOrbit} in Lemma M;
quotienting by a larger conjugacy
group has not lost multiplicities.
\end{proof}

\paragraph{PX9: lifting derivations and the exact one-dimensional tangent bound.}
\label{claim:PX9}
Assume \texttt{Fin0}, \texttt{NoEnd} and
\texttt{AllIrr RigidAt}.  If $x:A\to K$ vanishes on
$I$, is two-sided, and $\delta$ is an $O$-linear
point derivation at $x$ killing $I$, then there is
$b\in K$ such that
\[
                   \delta(f)=b n x(f)
                       \quad(f\in A_n).
 \tag{PX.23}
\]
In particular the exact predicate $\mathtt{M4}$ holds.

\begin{proof}
Package the derivation as
$X=x+\epsilon\delta:A\to K[\epsilon]$.
PX4 first gives a lift $\xi_0:\mathcal S\to K$
of $x$.  PX7 makes its representation
irreducible.  Hence there are entry functions
$u\in\mathcal S_{(1,0)}$ and
$v\in\mathcal S_{(-1,0)}$ with
$\xi_0(u)\xi_0(v)\ne0$ (use entries
$(0,1)$ and $(1,0)$).  PX4 gives
$a_0\in A_0$ with $\varphi(a_0)=uv$.
The element $c=X(a_0)$ is a unit in
$K[\epsilon]$.

This supplies an explicit extension of
$X$ to $\mathcal S$.  After inverting
$uv$, the first-weight grading has
an invertible element $u$ of degree
$1$ and its degree-zero subring is
$\varphi(A)[(uv)^{-1}]$.  Every element
of first degree $a$ is uniquely $u^a$
times an element of that subring.
Send $u$ to $1$ and extend the map
$X$ on the degree-zero subring.
In nonlocalized notation, multiply
a homogeneous $f$ of first degree
$a$ by $u^{(-a)_+}v^{a_+}$, lift
this balanced element to $A$, apply
$X$, and divide by $c^{a_+}$.
The kernel assertion of PX4 proves
independence of the lift; balancing
two factors proves multiplicativity.
This gives
$\Xi:\mathcal S\to K[\epsilon]$
with $\Xi\varphi=X$.  Its reduction
need not equal the initially chosen
$\xi_0$; its restriction to $A$ is
$x$, which is exactly what is needed.

Scale the reduction of $\Xi$ to an
actual point using PX5, and use the
same constant torus element on $\Xi$.
The resulting representation over
$K[\epsilon]$ is itself an actual
point.  We include the reason for
this sometimes delicate passage.
The reduced entries satisfy (PX.15).
The span assertion of PX5 implies
that all reduced elements of
$\mathcal S$ have norm at most $1$,
and reduced formal variables have
norm strictly less than $1$; to
check the latter subtract their
original scalar residues and use
the same bound on entry polynomials.
In a dual-number algebra, an element
with integral first component is
integral over $O$: its nilpotent
part is integral, and integrality
is closed under addition.  Adjoin
the finitely many variable values,
the finite module of weight-zero
values, and sufficiently large
powers of the variables divided
by a uniformizer, just as in PX5.
The resulting finite $O$-algebra
is complete, including its nilpotent
ideal; the variables are in the
radical of the uniformizer ideal.
The identical finite-monomial and
truncation argument extends $\Xi$
to $S\to K[\epsilon]$.  Thus the
original universal \texttt{IsPt}
lift is an actual witness for
$D.\mathtt{PtIn}$, as required by
the rigidity hypothesis.

Write this reduced actual point as
$\rho$; it is irreducible by PX7.
Rigidity gives the derivative
$[Y,\rho]$ with
$Y\in\mathfrak{sp}_4(K)$.
The differentiated block normalization
at $\phi$ and the maximum-norm
argument of PX8 force $Y$ to be
block diagonal.  Differentiating
the two zero corners of (PX.16)
then forces $Y_{02}=Y_{20}=0$,
because $H_{00}-H_{22}\ne0$.
The symplectic Lie-algebra equation
kills $Y_{13},Y_{31}$ as well.
Thus $Y$ is diagonal, and there
are $a,b\in K$ with
\[
 [Y,\rho(g)]_{ij}
   =(a(\lambda_i-\lambda_j)
                      +b(\mu_i-\mu_j))\rho(g)_{ij}.
 \tag{PX.24}
\]
This is the infinitesimal torus
$(1+\epsilon a,1+\epsilon b)$.
Its action on the constant point
has the same entries as the
dual-number point.  PX5 gives
equality of the maps on all of
$\mathcal S$.  On $A_n$, the
first-weight contribution is zero.
Undo the earlier constant scaling;
its nonzero factor cancels from
both sides, giving (PX.23).

For two derivations with parameters
$b,b'$, choose coefficients
$(b',-b)$ unless $b=0$, in which
case choose $(1,0)$.  They are
not both zero and their linear
combination vanishes on every
homogeneous piece by (PX.23),
hence on $A$.  This proves the
literal $\mathtt{M4}$, including
the actual-point restriction in
its use of rigidity.
\end{proof}

\paragraph{PX10: the trace element, orbit cover and actual Selmer inequality.}
\label{claim:PX10}
Under the same three arithmetic assertions, let
$g\in G$ satisfy
\[
 \operatorname{tr}\rho(g)-\operatorname{tr}\rho_0(g)\ne0
       \quad\text{for every irreducible actual point }\rho.
 \tag{PX.25}
\]
Then the literal source inequality $D.\mathtt{Ineq}(g)$
holds.  In particular its Selmer length is finite and
there is a finite family $\mathcal X$ of irreducible
actual representations, with no two distinct members
conjugate, such that
\[
 \operatorname{nrm}_O(D.\mathtt{sha})\log\ell
 \ \le\ 2\sum_{\rho\in\mathcal X}
 -\log\left|\operatorname{tr}\rho(g)
                         -\operatorname{tr}\rho_0(g)\right|.
 \tag{PX.26}
\]
This closes the entire internal model-to-comparison
part of SP's $G2$, with exactly \texttt{Fin0},
\texttt{NoEnd} and \texttt{AllIrr RigidAt} retained
as arithmetic input predicates.

\begin{proof}
Choose weight-zero lifts $f_i\in R_0$ of the
four diagonal entries of $\rho^u(g)$.  Set
\[
           a_g=\sum_{i=0}^3 f_i
                       -\operatorname{tr}\rho_0(g)\in A_0.
 \tag{PX.27}
\]
The model's origin identity gives $x_0(a_g)=0$;
at every point of $\mathcal S$ its value is
exactly the trace difference.  Torus
conjugation leaves this value unchanged.
PX7 and (PX.25) therefore prove $\mathtt{M5}$.

All the other hypotheses of the full
intersection theorem SP6 are now proved
on this same $A$: PX1 gives the weighted
ring, PX2--PX3 give $\mathtt{M1}$, PX5
gives $\mathtt{M2}$, PX7 gives
$\mathtt{M3}$, and PX9 gives
$\mathtt{M4}$.  Use the original
normalized valuation
\[
 v(z)=\frac{v_O(\ell)}{\log\ell}(-\log|z|),
                  \qquad |\ell|=\ell^{-1}.
 \tag{PX.28}
\]
Multiplicativity, the ultrametric inequality
and the DVR normalization check precisely
$\mathtt{IsNVal}$; no coefficient
extension is introduced by this step.
SP6 supplies a finite cover of all
two-sided orbits and the bound
\[
 \operatorname{length}_O C_{-1}(F)
                       \le2\sum_{x\in T}v(x(a_g)).
 \tag{PX.29}
\]
Choose such a cover of minimum cardinality.
It contains only two-sided points and
contains no two points in the same orbit,
since one of those could be deleted.
Lift and scale each point by PX4--PX5.
They are irreducible by PX7, and PX8
shows that distinct chosen points give
nonconjugate representations.  Thus
their trace values give exactly the
sum in (PX.26), without a multiplicity
loss or an enlargement of the set.

SP7 supplies the actual normalized
Selmer-cocycle injection into the
cotangent quotient on $M$, using
the model's original derivation
field and RV3's finiteness.  Together
with SP1 it bounds the original
Selmer length by the left side of
(PX.29), and proves that length is
finite before its real normalization.
Divide (PX.29) by $v_O(\ell)$ and
multiply by $\log\ell$ to obtain
(PX.26).  The factor is exactly
$2$.  Every representation, trace,
origin, equation, lattice and
coefficient normalization is the
one attached to the original $D$.
\end{proof}

\subsection{What the completed comparison leaves to the arithmetic count}

MX constructs the actual model from the local and
global installation in DX; AX constructs the original
$\mathtt{Good}$/$\mathtt{OfCase}$ datum and its exact
level and ordinary specializations.  Consequently
PX10 removes the internal weighted algebra, point
lifting, reducibility classification, tangent lifting
and orbit comparison from SP's remaining $G2$ input.
To use the paired-count conclusion SP10, the remaining
arithmetic theorem must still supply
$\mathtt{Fin0}$, $\mathtt{NoEnd}$, actual irreducible
rigidity and the common admissible representation
class with its uniform finite catalogue.  The image
theorem supplying the first two assertions is the
same one used in AX's final admissibility join.
This is a precise sufficient cut, not a proof of
the image theorem or a claim that an arbitrary
finite set of representations is an adequate
catalogue.  Once these inputs are supplied on
the specified objects, SP10 gives actual
$\mathtt{CountII}$ and $\mathtt{CountIII}$, and
HV10 already supplies their complete downstream
place in the original-field cohomological window
assembly.  The other actual counts and the
common-family, gap, theta, denominator and
metric/spectral branches keep their separate
roles in the full RH/Dirichlet-GRH reconstruction.
